\chapter{Deep Learning}

\section{Activation Functions}

\subsection{Sigmoid Activation Function}

The sigmoid function $\sigma(x)$ is defined as
\be
\sigma(x) = \frac{1}{1 + \exp(-x)},
\ee
shown in Fig.~\ref{fig:sigmoid_activation}

\begin{figure}[htb]
  \begin{center}
    \begin{overpic}[width=\textwidth]{fig/sigmoid_activation.pdf}
    %\put(17, 43.5){$\cA$}
    \end{overpic}
   \end{center}
  \caption{Sigmoid activation function.}
  \label{fig:sigmoid_activation}
\end{figure}


\begin{Hresult}
The derivative of the sigmoid activation function
$\sigma'(x)$ is 
\be
\sigma'(x) = \sigma \left( 1 - \sigma \right).
\ee
\end{Hresult}

\begin{Hproof}
\begin{align}
\sigma'(x)  
& = 
\frac{d}{dx} \left((1 + \exp(-x))^{-1}  \right) \\
& = 
(-1) (1 + \exp(-x))^{-2} \exp(-x) (-1) \\
& = 
\frac{\exp(-x)}{(1 + \exp(-x))^2} \\
& = 
\frac{(1 + \exp(-x)) - 1}{(1 + \exp(-x))^2} \\
& = 
\frac{1 + \exp(-x)}{(1 + \exp(-x))^2} - 
\left(\frac{1}{1 + \exp(-x)}\right)^2 \\
& = 
\sigma - \sigma^2 \\
& = 
\sigma (1 - \sigma)
\end{align}
\end{Hproof}


\footnote{See nice tutorials by Ryan Harris at
https://www.youtube.com/user/nqramjets}



